\documentclass[a4paper,11pt]{article}

\usepackage[a4paper,margin=2cm]{geometry}
\usepackage[T1]{fontenc}
\usepackage[utf8]{inputenc}
\usepackage[ngerman,english]{babel}
\usepackage{lmodern}
\usepackage{microtype}
\usepackage{xcolor}
\usepackage{colortbl}
\usepackage{array}
\usepackage{parskip}
\usepackage{enumitem}
\usepackage[most]{tcolorbox}
\usepackage{titlesec}
\usepackage[hidelinks]{hyperref}

\definecolor{HeaderBlueDark}{RGB}{70,110,170}
\definecolor{RowBlueLight}{RGB}{235,243,255}
\definecolor{RowBlueDark}{RGB}{220,233,250}
\definecolor{SoftGray}{RGB}{90,90,90}

\setlength{\parindent}{0pt}
\setlist[itemize]{leftmargin=1.4em,itemsep=0.35em,topsep=0.35em}
\linespread{1.06}

\titleformat{\section}[block]
{\Large\bfseries\color{HeaderBlueDark}}
{}{0pt}{}
\titlespacing{\section}{0pt}{1.3em}{0.8em}

\titleformat{\subsection}[block]
{\large\bfseries\color{HeaderBlueDark}}
{}{0pt}{}
\titlespacing{\subsection}{0pt}{1.0em}{0.6em}

\newtcolorbox{exbox}{enhanced,breakable,colback=RowBlueLight,colframe=RowBlueDark,boxrule=0.4pt,arc=1.5mm,left=1.2mm,right=1.2mm,top=1mm,bottom=1mm,title={Beispiele \textnormal{(Examples)}},fonttitle=\bfseries,colbacktitle=RowBlueDark,coltitle=black}

\NewDocumentEnvironment{grammarentry}{mm+b}{%
    \begin{tcolorbox}[enhanced,breakable,colback=white,colframe=HeaderBlueDark,boxrule=0.6pt,arc=1.5mm,left=1.5mm,right=1.5mm,top=1mm,bottom=1mm,colbacktitle=HeaderBlueDark,coltitle=white,fonttitle=\bfseries,title={#1\hfill\normalfont\footnotesize #2}]
    #3
    \end{tcolorbox}
}{}

\newcommand{\GermanText}[1]{\textbf{Deutsch: }#1\par}
\newcommand{\EnglishText}[1]{{\small\color{SoftGray}\textbf{English: }#1\par}}
\newcommand{\ExampleItem}[2]{\item \textbf{#1}\\{\small\color{SoftGray}#2}}

\title{Schwache, starke und gemischte Verben}
\date{}

\begin{document}
    \maketitle
    \tableofcontents
    \newpage
    
    
    \section{Grundidee}
        
        \begin{grammarentry}{Was ist wichtig?}{What is important?}
            \GermanText{Im Deutschen gibt es drei wichtige Gruppen von Verben: schwache Verben, starke Verben und gemischte Verben.}
            \EnglishText{In German there are three important groups of verbs: weak verbs, strong verbs, and mixed verbs.}
            
            \GermanText{Diese Gruppen sind besonders wichtig für das Präteritum und das Perfekt.}
            \EnglishText{These groups are especially important for the simple past and the present perfect.}
            
            \GermanText{Man muss bei jedem Verb wissen, ob es schwach, stark oder gemischt ist.}
            \EnglishText{For every verb, one must know whether it is weak, strong, or mixed.}
        \end{grammarentry}
    
    
    \section{Die drei Verbgruppen}
        
        \renewcommand{\arraystretch}{1.35}
        
        \begin{tabular}{>{\raggedright\arraybackslash}p{3.2cm} >{\raggedright\arraybackslash}p{4.2cm} >{\raggedright\arraybackslash}p{4.2cm} >{\raggedright\arraybackslash}p{3.2cm}}
            \rowcolor{HeaderBlueDark}
            \textcolor{white}{\textbf{Gruppe}} &
            \textcolor{white}{\textbf{Präteritum}} &
            \textcolor{white}{\textbf{Perfekt / Partizip II}} &
            \textcolor{white}{\textbf{Beispiel}} \\
            
            \rowcolor{RowBlueLight}
            schwache Verben &
            Endung: -te &
            ge- + Stamm + -t &
            machen, machte, gemacht \\
            
            \rowcolor{RowBlueDark}
            starke Verben &
            Vokalwechsel, keine -te-Endung &
            ge- + veränderter Stamm + -en &
            gehen, ging, gegangen \\
            
            \rowcolor{RowBlueLight}
            gemischte Verben &
            Vokalwechsel + -te &
            ge- + veränderter Stamm + -t &
            bringen, brachte, gebracht \\
        \end{tabular}
    
    
    \section{Schwache Verben}
        
        \begin{grammarentry}{Schwache Verben}{Weak verbs}
            \GermanText{Schwache Verben sind regelmäßig. Sie ändern den Stammvokal im Präteritum und im Partizip II normalerweise nicht.}
            \EnglishText{Weak verbs are regular. They normally do not change the stem vowel in the simple past and in the past participle.}
            
            \GermanText{Das Präteritum bildet man meistens mit -te. Das Partizip II bildet man meistens mit ge- und -t.}
            \EnglishText{The simple past is usually formed with -te. The past participle is usually formed with ge- and -t.}
            
            \begin{exbox}
                \begin{itemize}
                    \ExampleItem{machen: ich machte, ich habe gemacht}{to do/make: I did/made, I have done/made}
                    \ExampleItem{lernen: ich lernte, ich habe gelernt}{to learn: I learned, I have learned}
                    \ExampleItem{kaufen: ich kaufte, ich habe gekauft}{to buy: I bought, I have bought}
                \end{itemize}
            \end{exbox}
        \end{grammarentry}
        
        \subsection{Schema für schwache Verben}
            
            \begin{tabular}{>{\raggedright\arraybackslash}p{4cm} >{\raggedright\arraybackslash}p{5cm} >{\raggedright\arraybackslash}p{5cm}}
                \rowcolor{HeaderBlueDark}
                \textcolor{white}{\textbf{Form}} &
                \textcolor{white}{\textbf{Schema}} &
                \textcolor{white}{\textbf{Beispiel}} \\
                
                \rowcolor{RowBlueLight}
                Präsens &
                Stamm + normale Endung &
                ich mache \\
                
                \rowcolor{RowBlueDark}
                Präteritum &
                Stamm + -te &
                ich machte \\
                
                \rowcolor{RowBlueLight}
                Perfekt &
                haben + ge- + Stamm + -t &
                ich habe gemacht \\
            \end{tabular}
    
    
    \section{Starke Verben}
        
        \begin{grammarentry}{Starke Verben}{Strong verbs}
            \GermanText{Starke Verben sind unregelmäßig. Sie ändern oft den Stammvokal.}
            \EnglishText{Strong verbs are irregular. They often change the stem vowel.}
            
            \GermanText{Im Präteritum haben sie normalerweise keine -te-Endung. Das Partizip II endet meistens auf -en.}
            \EnglishText{In the simple past they normally do not have the ending -te. The past participle usually ends in -en.}
            
            \begin{exbox}
                \begin{itemize}
                    \ExampleItem{gehen: ich ging, ich bin gegangen}{to go: I went, I have gone}
                    \ExampleItem{finden: ich fand, ich habe gefunden}{to find: I found, I have found}
                    \ExampleItem{schreiben: ich schrieb, ich habe geschrieben}{to write: I wrote, I have written}
                \end{itemize}
            \end{exbox}
        \end{grammarentry}
        
        \subsection{Schema für starke Verben}
            
            \begin{tabular}{>{\raggedright\arraybackslash}p{4cm} >{\raggedright\arraybackslash}p{5cm} >{\raggedright\arraybackslash}p{5cm}}
                \rowcolor{HeaderBlueDark}
                \textcolor{white}{\textbf{Form}} &
                \textcolor{white}{\textbf{Schema}} &
                \textcolor{white}{\textbf{Beispiel}} \\
                
                \rowcolor{RowBlueLight}
                Präsens &
                manchmal Vokalwechsel in du/er/sie/es &
                du fährst \\
                
                \rowcolor{RowBlueDark}
                Präteritum &
                Stammvokal ändert sich &
                ich fuhr \\
                
                \rowcolor{RowBlueLight}
                Perfekt &
                haben/sein + ge- + veränderter Stamm + -en &
                ich bin gefahren \\
            \end{tabular}
    
    
    \section{Gemischte Verben}
        
        \begin{grammarentry}{Gemischte Verben}{Mixed verbs}
            \GermanText{Gemischte Verben haben Eigenschaften von schwachen und starken Verben.}
            \EnglishText{Mixed verbs have features of both weak and strong verbs.}
            
            \GermanText{Sie haben einen Vokalwechsel wie starke Verben, aber sie bekommen im Präteritum die Endung -te wie schwache Verben.}
            \EnglishText{They have a vowel change like strong verbs, but in the simple past they take the ending -te like weak verbs.}
            
            \GermanText{Das Partizip II endet normalerweise auf -t, nicht auf -en.}
            \EnglishText{The past participle normally ends in -t, not in -en.}
            
            \begin{exbox}
                \begin{itemize}
                    \ExampleItem{bringen: ich brachte, ich habe gebracht}{to bring: I brought, I have brought}
                    \ExampleItem{denken: ich dachte, ich habe gedacht}{to think: I thought, I have thought}
                    \ExampleItem{kennen: ich kannte, ich habe gekannt}{to know: I knew, I have known}
                \end{itemize}
            \end{exbox}
        \end{grammarentry}
        
        \subsection{Schema für gemischte Verben}
            
            \begin{tabular}{>{\raggedright\arraybackslash}p{4cm} >{\raggedright\arraybackslash}p{5cm} >{\raggedright\arraybackslash}p{5cm}}
                \rowcolor{HeaderBlueDark}
                \textcolor{white}{\textbf{Form}} &
                \textcolor{white}{\textbf{Schema}} &
                \textcolor{white}{\textbf{Beispiel}} \\
                
                \rowcolor{RowBlueLight}
                Präsens &
                normale Präsensform &
                ich bringe \\
                
                \rowcolor{RowBlueDark}
                Präteritum &
                Vokalwechsel + -te &
                ich brachte \\
                
                \rowcolor{RowBlueLight}
                Perfekt &
                haben + ge- + veränderter Stamm + -t &
                ich habe gebracht \\
            \end{tabular}
    
    
    \section{Sehr wichtige gemischte Verben}
        
        \begin{tabular}{>{\raggedright\arraybackslash}p{3.5cm} >{\raggedright\arraybackslash}p{3.5cm} >{\raggedright\arraybackslash}p{4cm} >{\raggedright\arraybackslash}p{3.5cm}}
            \rowcolor{HeaderBlueDark}
            \textcolor{white}{\textbf{Infinitiv}} &
            \textcolor{white}{\textbf{Präteritum}} &
            \textcolor{white}{\textbf{Perfekt}} &
            \textcolor{white}{\textbf{Englisch}} \\
            
            \rowcolor{RowBlueLight}
            bringen &
            brachte &
            hat gebracht &
            to bring \\
            
            \rowcolor{RowBlueDark}
            denken &
            dachte &
            hat gedacht &
            to think \\
            
            \rowcolor{RowBlueLight}
            kennen &
            kannte &
            hat gekannt &
            to know \\
            
            \rowcolor{RowBlueDark}
            nennen &
            nannte &
            hat genannt &
            to name/call \\
            
            \rowcolor{RowBlueLight}
            rennen &
            rannte &
            ist gerannt &
            to run \\
            
            \rowcolor{RowBlueDark}
            brennen &
            brannte &
            hat gebrannt &
            to burn \\
            
            \rowcolor{RowBlueLight}
            wissen &
            wusste &
            hat gewusst &
            to know \\
        \end{tabular}
    
    
    \section{Wichtige Regel für das Perfekt}
        
        \begin{grammarentry}{haben oder sein?}{haben or sein?}
            \GermanText{Die meisten Verben bilden das Perfekt mit haben.}
            \EnglishText{Most verbs form the present perfect with haben.}
            
            \GermanText{Viele Verben der Bewegung oder Zustandsänderung bilden das Perfekt mit sein.}
            \EnglishText{Many verbs of movement or change of state form the present perfect with sein.}
            
            \begin{exbox}
                \begin{itemize}
                    \ExampleItem{Ich habe gelernt.}{I have learned.}
                    \ExampleItem{Ich bin gegangen.}{I have gone.}
                    \ExampleItem{Ich bin gerannt.}{I have run.}
                \end{itemize}
            \end{exbox}
        \end{grammarentry}
    
    
    \section{Praktische Lernregel}
        
        \begin{grammarentry}{Was muss man lernen?}{What does one have to learn?}
            \GermanText{Bei schwachen Verben reicht meistens der Infinitiv, weil die Formen regelmäßig sind.}
            \EnglishText{With weak verbs, the infinitive is usually enough because the forms are regular.}
            
            \GermanText{Bei starken und gemischten Verben sollte man immer drei Formen lernen: Infinitiv, Präteritum und Perfekt.}
            \EnglishText{With strong and mixed verbs, one should always learn three forms: infinitive, simple past, and present perfect.}
            
            \begin{exbox}
                \begin{itemize}
                    \ExampleItem{gehen -- ging -- ist gegangen}{to go -- went -- has gone}
                    \ExampleItem{finden -- fand -- hat gefunden}{to find -- found -- has found}
                    \ExampleItem{bringen -- brachte -- hat gebracht}{to bring -- brought -- has brought}
                \end{itemize}
            \end{exbox}
        \end{grammarentry}
    
    
    \section{Kurz zusammengefasst}
        
        \begin{tabular}{>{\raggedright\arraybackslash}p{4cm} >{\raggedright\arraybackslash}p{5cm} >{\raggedright\arraybackslash}p{5cm}}
            \rowcolor{HeaderBlueDark}
            \textcolor{white}{\textbf{Gruppe}} &
            \textcolor{white}{\textbf{Merkmal}} &
            \textcolor{white}{\textbf{Beispiel}} \\
            
            \rowcolor{RowBlueLight}
            schwach &
            kein Vokalwechsel, Präteritum mit -te, Partizip II mit -t &
            machen -- machte -- gemacht \\
            
            \rowcolor{RowBlueDark}
            stark &
            Vokalwechsel, Präteritum ohne -te, Partizip II meistens mit -en &
            gehen -- ging -- gegangen \\
            
            \rowcolor{RowBlueLight}
            gemischt &
            Vokalwechsel + -te, Partizip II mit -t &
            bringen -- brachte -- gebracht \\
        \end{tabular}

\end{document}